\section{Introduction}
Plusieurs techniques de rendu n'utilisant pas une représentation par des triangles sont devenues populaires récemment, que ce soit les NeRF \cite{mildenhall2021nerf} ou le Gaussian splatting \cite{kerbl2023gaussian}, ces techniques gagnent en popularité puisqu'elles permettent de représenter la scène sans avoir recours à des géométries maillées en triangle. De plus, ces algorithmes s'inscrivent dans un courant tentant de rendre toutes les opérations de rendu différentiables \cite{kato2020differentiable}. Cela permet d'entraîner des réseaux de neurones sur la représentation de la scène. Récemment, une nouvelle technique a été publiée par l'équipe de \textcite{seyb24from} ; cette méthode représente les éléments de la scène comme des volumes encodant un processus gaussien (GPIS), elle est aussi différentiable.

Lors d'un rendu maillé traditionnel, les triangles sont rendus par tracé de rayon, alors que les volumes et la fumée utilisent des algorithmes de tracé de rayon volumétrique comme le \textit{raymarching}, nécessitant ainsi le support de plusieurs algorithmes de tracé de rayon différents dans un seul moteur de rendu. L'approche GPIS permet de représenter les surfaces et les volumes avec le même algorithme, permettant ainsi de faire le rendu d'une scène avec un seul algorithme.

En ce sens, il est intéressant d'évaluer la faisabilité du rendu de tous les éléments d'une scène par l'algorithme GPIS. D'autre part, un élément omniprésent dans la littérature sur les techniques de rendu est la chevelure, qui représente traditionnellement un défi par la quantité extrêmement élevée de cheveux se trouvant sur la tête d'un personnage. Dans ce projet de recherche, je vais appliquer et étendre la technique de rendu par GPIS pour permettre la représentation des cheveux.

Dans ce rapport, je présenterai le contexte nécessaire pour comprendre la technique de rendu par GPIS ainsi que l'état de l'art sur les techniques de représentation des cheveux. Ensuite, je présenterai le modèle développé dans ce projet de recherche et finalement, les rendus obtenus seront analysés et discutés.
